% ==========================================================
% HE2002 Macroeconomics II — Chapter 1: The Goods Market
% Source materials: :contentReference[oaicite:0]{index=0}
% (This is a chapter file intended to be \input into the main template.)
% ==========================================================

\section{The Goods Market}

\subsection{GDP and the Composition of Demand}

\begin{definition}{GDP (expenditure approach) and demand components}{def:gdp_components}
Gross Domestic Product (GDP) is commonly measured using the expenditure approach:
\emph{total spending on domestically produced final goods and services}.
Key components are:
consumption $C$, investment $I$, government purchases $G$, exports $X$, and imports $IM$.
\end{definition}

\begin{keyformula}[title={Key Formula: Total demand for goods}]
\[
Z \equiv C + I + G + X - IM,
\qquad
NX \equiv X-IM.
\]
\end{keyformula}

\noindent
In these notes we focus on a \textbf{closed economy} (the baseline for HE2002 short-run analysis), so $X=IM=0$ and
\[
Z \equiv C + I + G.
\]
We also \textbf{ignore inventory investment} in the model; nevertheless, inventory adjustments provide a useful
economic mechanism for how output adjusts when production differs from desired spending (see below).

\subsection{Model Setup: Behavioural Equations and Policy Variables}

\begin{definition}{Disposable income and net taxes}{def:disposable_income}
Disposable income is income available to households after taxes net of transfers:
\[
Y_D \equiv Y - T,
\]
where $Y$ is income/output and $T$ denotes taxes minus government transfers (net taxes).
\end{definition}

\begin{definition}{Consumption function (baseline Keynesian specification)}{def:consumption_function}
Consumption is assumed to depend positively on disposable income:
\[
C = C(Y_D), \qquad \frac{dC}{dY_D} > 0.
\]
A standard linear specification is
\[
C = c_0 + c_1 Y_D,
\quad 0 < c_1 < 1,
\]
where $c_1$ is the propensity to consume and $c_0$ is autonomous consumption.
\end{definition}

\begin{keyformula}[title={Key Formula: Consumption written in terms of $Y$ and $T$}]
Using $Y_D=Y-T$,
\[
C = c_0 + c_1(Y-T).
\]
\end{keyformula}

\begin{definition}{Endogenous vs.\ exogenous variables}{def:endo_exo}
\textbf{Endogenous} variables are determined within the model (e.g.\ equilibrium output $Y$).
\textbf{Exogenous} variables are taken as given (e.g.\ policy instruments or parameters).
A bar (e.g.\ $\bar I$) indicates an exogenous constant.
\end{definition}

\noindent
For the baseline goods-market model:
\begin{itemize}
\item Investment is taken as exogenous: $I=\bar I$ (this assumption is later relaxed in IS--LM analysis).
\item Fiscal policy variables $G$ and $T$ are taken as exogenous choices of the government.
\item We ignore the interest-rate effect on consumption in this chapter.
\end{itemize}

\subsection{Goods-Market Equilibrium and Equilibrium Output}

\begin{keyformula}[title={Key Formula: Planned expenditure in a closed economy}]
With $I=\bar I$,
\[
Z = C + \bar I + G
  = \bigl[c_0 + \bar I + G - c_1T\bigr] + c_1Y.
\]
\end{keyformula}

\begin{theorem}{Equilibrium output in the baseline goods-market model}{equilibrium_output}
Goods-market equilibrium requires output to equal demand:
\[
Y = Z.
\]
Under the linear consumption function and $I=\bar I$, equilibrium output is
\[
Y
= \frac{1}{1-c_1}\Bigl[c_0 + \bar I + G - c_1T\Bigr],
\qquad 0<c_1<1.
\]
\end{theorem}

\noindent
\textbf{Derivation.} Start from the equilibrium condition $Y=Z$ and substitute for $Z$:
\[
Y = c_0 + c_1(Y-T) + \bar I + G
  = c_0 + c_1Y - c_1T + \bar I + G.
\]
Collect $Y$ terms on the left:
\[
(1-c_1)Y = c_0 + \bar I + G - c_1T.
\]
Divide by $(1-c_1)$ (positive since $0<c_1<1$) to obtain the stated expression.

\begin{definition}{Autonomous spending}{def:autonomous_spending}
In the baseline model, define \emph{autonomous spending} as the component of planned expenditure not
depending on output:
\[
A \equiv c_0 + \bar I + G - c_1T.
\]
Then equilibrium output can be written compactly as $Y=\dfrac{1}{1-c_1}A$.
\end{definition}

\begin{commonmistake}
Do not confuse:
\begin{itemize}
\item the \textbf{accounting identity} defining demand ($Z \equiv C+I+G$ in a closed economy), and
\item the \textbf{equilibrium condition} ($Y=Z$), which pins down equilibrium output.
\end{itemize}
The identity holds by definition; the equilibrium condition is an economic restriction that holds when the goods market clears.
\end{commonmistake}

\subsection{Adjustment Mechanism and the Keynesian Cross Intuition}

\noindent
Even though we ignore inventory investment in the model, unintended inventories provide intuition for why output
moves toward equilibrium:
\begin{itemize}
\item If $Y<Z$, desired spending exceeds production; inventories fall unexpectedly and firms raise production.
\item If $Y>Z$, production exceeds desired spending; inventories accumulate unexpectedly and firms cut production.
\end{itemize}

\noindent
Graphically, the Keynesian cross plots $Z$ against $Y$ together with the $45^\circ$ line (where $Z=Y$).
The intersection determines equilibrium output.

\begin{center}
\begin{tikzpicture}[scale=0.95]
  \draw[->] (0,0) -- (8,0) node[below] {$Y$};
  \draw[->] (0,0) -- (0,6) node[left] {$Z$};

  % 45-degree line
  \draw (0,0) -- (7.2,5.4) node[above right] {$Z=Y$};

  % Demand line Z = A + c1 Y (illustrative)
  \draw (0,1.2) -- (7.2,4.8) node[above right] {$Z = A + c_1Y$};

  % Mark an equilibrium point (illustrative)
  \draw[dashed] (4.8,0) -- (4.8,3.6);
  \draw[dashed] (0,3.6) -- (4.8,3.6);
  \fill (4.8,3.6) circle (2pt) node[above left] {$Y^\ast$};
\end{tikzpicture}
\end{center}

\subsection{Fiscal Multipliers in the Baseline Model}

\begin{keyformula}[title={Key Formula: Multipliers (holding other exogenous variables fixed)}]
From
\(
Y=\dfrac{1}{1-c_1}\bigl[c_0+\bar I+G-c_1T\bigr],
\)
\[
\frac{\partial Y}{\partial G} = \frac{1}{1-c_1},
\qquad
\frac{\partial Y}{\partial T} = -\frac{c_1}{1-c_1}.
\]
\end{keyformula}

\noindent
\textbf{Interpretation.}
Government spending enters demand one-for-one, while taxes affect demand only through consumption.
A one-unit increase in $T$ reduces disposable income by one unit, but consumption falls only by $c_1$ units.

\begin{keyformula}[title={Key Formula: Geometric-series intuition for the multiplier}]
If autonomous spending increases by $\Delta A$, then output rises in successive rounds:
\[
\Delta Y
= \Delta A + c_1\Delta A + c_1^2\Delta A + \cdots
= \frac{1}{1-c_1}\Delta A.
\]
\end{keyformula}

\begin{examplebox}[title={Worked Example (Tutorial 1, Question 1): Equilibrium GDP, disposable income, consumption}]
\textbf{Given} (billions of euros):
\[
C = 480 + 0.5Y_D,\quad I=110,\quad T=70,\quad G=250.
\]
Solve for:
\begin{enumerate}[label=(\alph*)]
\item equilibrium output $Y$,
\item disposable income $Y_D$,
\item consumption $C$.
\end{enumerate}

\textbf{Solution.}
Disposable income is
\[
Y_D = Y - T = Y - 70.
\]
Substitute into the consumption function:
\[
C = 480 + 0.5(Y-70)
  = 480 + 0.5Y - 35
  = 445 + 0.5Y.
\]
Total demand (closed economy) is
\[
Z = C + I + G
  = (445 + 0.5Y) + 110 + 250
  = 805 + 0.5Y.
\]
Equilibrium requires $Y=Z$:
\[
Y = 805 + 0.5Y
\;\Rightarrow\;
0.5Y = 805
\;\Rightarrow\;
Y = 1610.
\]
Then
\[
Y_D = Y - 70 = 1610 - 70 = 1540,
\qquad
C = 480 + 0.5\cdot 1540 = 480 + 770 = 1250.
\]
\end{examplebox}


\newpage
\subsection{Balanced Budget Multiplier}

\begin{examplebox}[title={Worked Example (Tutorial 1, Question 3): Balanced budget multiplier}]
Start from the equilibrium expression
\[
Y = \frac{1}{1-c_1}\bigl[c_0 + \bar I + G - c_1T\bigr],
\qquad 0<c_1<1.
\]
\begin{enumerate}[label=(\alph*)]
\item By how much does $Y$ increase when $G$ increases by one unit?
\item By how much does $Y$ change when $T$ increases by one unit?
\item Why are the answers to (a) and (b) different?
\item Suppose initially $G=T$ and policy changes keep the budget balanced: $\Delta G=\Delta T=1$.
What is $\Delta Y$? Are balanced-budget changes neutral?
\item How does the value of $c_1$ affect your answer to (d)?
\end{enumerate}

\textbf{Solution.}
\begin{enumerate}[label=(\alph*)]
\item Differentiate with respect to $G$:
\[
\frac{\partial Y}{\partial G} = \frac{1}{1-c_1}.
\]
So a one-unit increase in $G$ raises $Y$ by $\dfrac{1}{1-c_1}$.

\item Differentiate with respect to $T$:
\[
\frac{\partial Y}{\partial T} = -\frac{c_1}{1-c_1}.
\]
So a one-unit increase in $T$ reduces $Y$ by $\dfrac{c_1}{1-c_1}$.

\item A change in $G$ shifts demand directly (one-for-one).
A change in $T$ affects demand only through consumption via disposable income, so only the fraction $c_1$ feeds into spending.

\item With $\Delta G=\Delta T=1$,
\[
\Delta Y
= \frac{1}{1-c_1}\Delta G - \frac{c_1}{1-c_1}\Delta T
= \frac{1}{1-c_1} - \frac{c_1}{1-c_1}
= 1.
\]
Thus a balanced-budget increase in $(G,T)$ is \textbf{not} neutral: output rises by one unit.

\item The balanced-budget multiplier is
\(
\Delta Y = 1
\)
for any $0<c_1<1$.
While $\dfrac{\partial Y}{\partial G}$ and $\dfrac{\partial Y}{\partial T}$ depend on $c_1$, their balanced-budget net effect does not.
\end{enumerate}
\end{examplebox}



\newpage
\subsection{Automatic Stabilizers: Income-Dependent Taxes}

\begin{definition}{Automatic stabilizers (in the goods-market model)}{def:automatic_stabilizers}
Taxes often increase with income. If net taxes follow
\(
T=t_0+t_1Y
\)
with $0<t_1<1$, then part of any increase in income is automatically taxed away,
dampening fluctuations in disposable income, consumption, and output.
\end{definition}

\begin{examplebox}[title={Worked Example (Tutorial 1, Question 4): Income-dependent taxes and the multiplier}]
Consider:
\[
C = c_0 + c_1Y_D,\qquad
T = t_0 + t_1Y,\qquad
Y_D = Y - T,
\]
with constant $G$ and $I=\bar I$, and $0<t_1<1$.
\begin{enumerate}[label=(\alph*)]
\item Solve for equilibrium output $Y$.
\item What is the multiplier on autonomous spending? Compare $t_1=0$ vs.\ $t_1>0$.
\item Explain why this tax system acts as an automatic stabilizer.
\end{enumerate}

\textbf{Solution.}
First compute disposable income:
\[
Y_D = Y - (t_0+t_1Y) = (1-t_1)Y - t_0.
\]
Then consumption is
\[
C = c_0 + c_1Y_D
  = c_0 + c_1\bigl[(1-t_1)Y - t_0\bigr]
  = (c_0 - c_1t_0) + c_1(1-t_1)Y.
\]

\begin{enumerate}[label=(\alph*)]
\item Goods-market equilibrium $Y=C+\bar I+G$ implies
\[
Y = (c_0 - c_1t_0) + c_1(1-t_1)Y + \bar I + G.
\]
Move the $Y$ term to the left:
\[
\bigl[1-c_1(1-t_1)\bigr]Y = (c_0 - c_1t_0) + \bar I + G.
\]
Hence
\[
Y = \frac{(c_0 - c_1t_0) + \bar I + G}{1-c_1(1-t_1)}.
\]

\item The multiplier on autonomous spending (the numerator term) is
\[
\frac{1}{1-c_1(1-t_1)}.
\]
If $t_1=0$, it reduces to $\dfrac{1}{1-c_1}$.
If $t_1>0$, then $1-c_1(1-t_1) > 1-c_1$, so the multiplier is smaller:
income-dependent taxes dampen the response of $Y$ to shocks.

\item When $Y$ rises, $T$ rises automatically, reducing $Y_D$ and hence moderating the increase in $C$.
When $Y$ falls, $T$ falls automatically, cushioning $Y_D$ and $C$.
This automatic offset stabilizes output without requiring discretionary policy changes.
\end{enumerate}
\end{examplebox}

\begin{policynote}
\textbf{Policy implication (short run).}
In the baseline goods-market framework, discretionary fiscal expansions (higher $G$ or lower $T$) increase equilibrium output through the multiplier.
With income-dependent taxes, the economy becomes less sensitive to shocks: automatic stabilizers reduce the effective multiplier, dampening both booms and recessions.
\end{policynote}

\subsection{Investment Equals Saving: An Equivalent Equilibrium Condition}

\begin{definition}{Private, public, and national saving}{def:saving_definitions}
Private saving is disposable income not consumed:
\[
S \equiv Y_D - C = Y - T - C.
\]
Public saving is the government budget surplus:
\[
T-G \quad \bigl(>0\ \text{surplus},\ <0\ \text{deficit}\bigr).
\]
National saving is the sum of private and public saving: $S + (T-G)$.
\end{definition}

\begin{keyformula}[title={Key Formula: Investment equals national saving (closed economy)}]
Starting from goods-market equilibrium $Y=C+I+G$,
\[
Y - T - C = I + G - T.
\]
Since $S \equiv Y-T-C$, this becomes
\[
S = I + G - T
\quad\Longleftrightarrow\quad
I = S + (T-G).
\]
\end{keyformula}

\begin{commonmistake}
In a closed economy with a government, the correct saving--investment relation is
\[
I = S + (T-G),
\]
not $I=S$.
Forgetting the public saving term $(T-G)$ leads to incorrect conclusions about how fiscal deficits relate to investment.
\end{commonmistake}

\noindent
\textbf{Deriving the same equilibrium output via saving.}
From $C=c_0+c_1(Y-T)$, private saving is
\[
S = Y - T - C
  = Y - T - c_0 - c_1(Y-T)
  = -c_0 + (1-c_1)(Y-T).
\]
Impose the equilibrium condition $I = S + (T-G)$ with $I=\bar I$:
\[
\bar I = \bigl[-c_0 + (1-c_1)(Y-T)\bigr] + (T-G).
\]
Expand and collect $Y$ terms:
\[
\bar I = -c_0 + (1-c_1)Y - (1-c_1)T + T - G
      = -c_0 + (1-c_1)Y + c_1T - G.
\]
Move constants to the right:
\[
(1-c_1)Y = c_0 + \bar I + G - c_1T
\quad\Rightarrow\quad
Y = \frac{1}{1-c_1}\bigl[c_0+\bar I+G-c_1T\bigr],
\]
which matches Theorem~\ref{thm:equilibrium_output}.

\begin{examtip}
High-frequency exam tasks for the goods-market chapter:
\begin{itemize}
\item Write down the model primitives and classify variables as endogenous vs.\ exogenous.
\item Derive equilibrium output algebraically from $Y=Z$; simplify carefully.
\item Compute $\frac{\partial Y}{\partial G}$ and $\frac{\partial Y}{\partial T}$; explain intuitively why the tax multiplier differs.
\item Solve for $Y$ under a balanced-budget change $\Delta G=\Delta T$ and interpret the result.
\item Use the saving framework to show $I=S+(T-G)$ and derive equilibrium output as a consistency check.
\item Explain (in words) how income-dependent taxes reduce the multiplier and why this is an automatic stabilizer.
\end{itemize}
\end{examtip}
